\documentclass[12pt, letterpaper]{article}

% ── Paquetes ──────────────────────────────────────────────────────────────────
\usepackage[utf8]{inputenc}
\usepackage[T1]{fontenc}
\usepackage{lmodern}
\usepackage[spanish]{babel}
\usepackage{geometry}
\usepackage{graphicx}
\usepackage{xcolor}
\usepackage{colortbl}
\usepackage{setspace}
\usepackage{microtype}
\usepackage{booktabs}
\usepackage{tabularx}
\usepackage{array}
\usepackage{longtable}
\usepackage{multirow}
\usepackage{fancyhdr}
\usepackage{hyperref}
\usepackage{parskip}
\usepackage{float}
\usepackage{caption}

\geometry{
  letterpaper,
  top=2.5cm, bottom=2.5cm,
  left=3cm,  right=3cm,
  headheight=15pt
}

% ── Colores ───────────────────────────────────────────────────────────────────
\definecolor{azulPUJ}{HTML}{1F3864}
\definecolor{azulMed}{HTML}{2E74B5}
\definecolor{azulClr}{HTML}{D6E4F7}
\definecolor{grisText}{HTML}{404040}
\definecolor{grisCl}{HTML}{F2F2F2}
\definecolor{lineColor}{HTML}{BFBFBF}
\definecolor{naranCl}{HTML}{FCE4D6}
\definecolor{verdeCl}{HTML}{E2EFDA}

% ── Columnas de tabla personalizadas ─────────────────────────────────────────
% P{w} = parbox justificado con ancho fijo
\newcolumntype{P}[1]{>{\raggedright\arraybackslash}p{#1}}
\newcolumntype{C}[1]{>{\centering\arraybackslash}p{#1}}

% ── Header / Footer ───────────────────────────────────────────────────────────
\pagestyle{fancy}
\fancyhf{}
\renewcommand{\headrulewidth}{0.4pt}
\renewcommand{\footrulewidth}{0.4pt}
\fancyhead[L]{\small\color{azulMed} Plataforma IoT de Monitoreo Industrial --- SAD}
\fancyhead[R]{\small\color{grisText} Pontificia Universidad Javeriana}
\fancyfoot[L]{\small\color{grisText} Arquitectura de Software · Taller 1 · Caso 9}
\fancyfoot[R]{\small\color{azulMed} \thepage}

% ── Hipervínculos ─────────────────────────────────────────────────────────────
\hypersetup{
  colorlinks=true,
  linkcolor=azulMed,
  urlcolor=azulMed,
  pdftitle={SAD -- Plataforma IoT de Monitoreo Industrial},
  pdfauthor={Mesa, Peñuela, Parra}
}

% ── Espaciado entre párrafos ──────────────────────────────────────────────────
\setlength{\parskip}{6pt}
\setlength{\parindent}{0pt}

% ─────────────────────────────────────────────────────────────────────────────
\begin{document}

% Evita anclas duplicadas de hyperref al reiniciar numeracion despues de portada.
\hypersetup{pageanchor=false}

% ══════════════════════════════════════════════════════════════════════════════
% PORTADA
% ══════════════════════════════════════════════════════════════════════════════
\begin{titlepage}
  \thispagestyle{empty}
  \centering

  
\includegraphics[width=4cm]{LogoPUJ.png}

  \vspace{0.4cm}
  {\color{lineColor}\rule{\linewidth}{0.4pt}}
  \vspace{0.3cm}

  {\fontsize{11}{13}\selectfont\color{azulPUJ}
    \textbf{PONTIFICIA UNIVERSIDAD JAVERIANA}
  }\\[0.15cm]
  {\fontsize{10}{12}\selectfont\color{grisText}
    Departamento de Ingeniería de Sistemas
  }\\[0.05cm]
  {\fontsize{10}{12}\selectfont\color{grisText}
    \textit{Arquitectura de Software}
  }

  \vspace{0.5cm}
  {\color{lineColor}\rule{\linewidth}{0.4pt}}

  \vfill

  {\fontsize{10}{12}\selectfont\color{grisText}
    \MakeUppercase{Documento de Arquitectura de Software}
  }\\[0.3cm]

  {\fontsize{20}{24}\selectfont\bfseries\color{azulPUJ}
    Plataforma IoT de Monitoreo\\[0.2cm]
    Industrial para Infraestructura\\[0.2cm]
    Energética
  }

  \vspace{0.6cm}

  {\fontsize{12}{14}\selectfont\color{grisText}
    Taller 1 \quad\textbullet\quad Caso 9
  }

  \vfill

  {\color{lineColor}\rule{\linewidth}{0.4pt}}
  \vspace{0.4cm}

  {\fontsize{10}{14}\selectfont\color{grisText}
    \begin{tabular}{c}
      Santiago Mesa \\[0.1cm]
      Camilo Peñuela \\[0.1cm]
    \end{tabular}
  }

  \vspace{0.5cm}
  {\color{lineColor}\rule{\linewidth}{0.4pt}}
  \vspace{0.3cm}

  {\fontsize{10}{12}\selectfont\color{grisText}
    Bogotá, Colombia \quad\textbullet\quad 11 de marzo de 2026
  }
\end{titlepage}

\hypersetup{pageanchor=true}

% Reinicia numeración desde la página de contenido
\pagenumbering{arabic}
\setcounter{page}{1}

% ══════════════════════════════════════════════════════════════════════════════
% SECCIÓN 1: REQUERIMIENTOS
% ══════════════════════════════════════════════════════════════════════════════
\section{Requerimientos}

El sistema consolida la telemetría de 120.000 dispositivos IoT distribuidos en subestaciones
eléctricas, plantas de generación y redes de transmisión nacionales. La plataforma procesa
aproximadamente 2,4 millones de eventos por minuto (3,4 mil millones diarios), detecta
anomalías críticas en menos de 3 segundos y garantiza una disponibilidad mínima del 99,95\%
anual. Los requerimientos se clasifican en tres categorías: proyecto, dominio y producto.

% ── 1.1 Requerimientos de Proyecto ───────────────────────────────────────────
\subsection{Requerimientos de Proyecto}

\begin{longtable}{
  C{1.5cm}
  P{2.8cm}
  P{5.2cm}
  P{4.2cm}
}
\rowcolor{azulPUJ}
\multicolumn{1}{>{\centering\arraybackslash\color{white}\bfseries}p{1.5cm}}{ID} &
\multicolumn{1}{>{\raggedright\arraybackslash\color{white}\bfseries}p{2.8cm}}{Nombre} &
\multicolumn{1}{>{\raggedright\arraybackslash\color{white}\bfseries}p{5.2cm}}{Descripción} &
\multicolumn{1}{>{\raggedright\arraybackslash\color{white}\bfseries}p{4.2cm}}{Justificación} \\
\endfirsthead
\rowcolor{azulPUJ}
\multicolumn{1}{>{\centering\arraybackslash\color{white}\bfseries}p{1.5cm}}{ID} &
\multicolumn{1}{>{\raggedright\arraybackslash\color{white}\bfseries}p{2.8cm}}{Nombre} &
\multicolumn{1}{>{\raggedright\arraybackslash\color{white}\bfseries}p{5.2cm}}{Descripción} &
\multicolumn{1}{>{\raggedright\arraybackslash\color{white}\bfseries}p{4.2cm}}{Justificación} \\
\endhead

\rowcolor{grisCl}
\textbf{RP-01} & Metodología Scrum &
  El equipo entrega incrementos funcionales cada dos semanas mediante sprints con revisiones
  y retrospectivas. &
  Permite responder con agilidad a cambios de requisitos en entornos industriales cambiantes. \\

\textbf{RP-02} & Arquitectura nativa en la nube &
  El diseño implementa microservicios en Docker orquestados con Kubernetes, aislando fallos
  operativos por módulo. &
  Evita el colapso total ante picos de alertas simultáneas; solo se reinicia el módulo afectado. \\

\rowcolor{grisCl}
\textbf{RP-03} & Tecnologías de código abierto &
  La plataforma integra InfluxDB, Apache Kafka y Grafana; evita licencias por nodo sensor. &
  Escala a cientos de miles de sensores sin incremento lineal de costos de licenciamiento. \\

\textbf{RP-04} & CI/CD automatizado &
  Pipeline de integración y despliegue continuo con pruebas automáticas previas a cada
  actualización productiva. &
  Elimina riesgo de fallos manuales en actualizaciones de un sistema crítico 24/7. \\

\rowcolor{grisCl}
\textbf{RP-05} & Pruebas de carga automatizadas &
  Simulación del tráfico equivalente a 120.000 sensores en entorno de pre-producción antes
  de cada release. &
  Detecta cuellos de botella antes de que ocurran en producción durante una emergencia real. \\

\textbf{RP-06} & Documentación como código &
  Arquitectura y decisiones técnicas documentadas en repositorios versionados junto al código
  fuente. &
  Garantiza trazabilidad de decisiones para auditorías regulatorias y nuevos integrantes del
  equipo. \\

\rowcolor{grisCl}
\textbf{RP-07} & Monitoreo de infraestructura propia &
  La plataforma mide continuamente CPU, memoria, latencia de red y throughput de los nodos
  de procesamiento. &
  Previene saturación silenciosa que retrasaría el procesamiento de alarmas críticas de la red
  eléctrica. \\

\end{longtable}

% ── 1.2 Requerimientos de Dominio ────────────────────────────────────────────
\subsection{Requerimientos de Dominio}

\begin{longtable}{
  C{1.5cm}
  P{2.8cm}
  P{5.2cm}
  P{4.2cm}
}
\rowcolor{azulPUJ}
\multicolumn{1}{>{\centering\arraybackslash\color{white}\bfseries}p{1.5cm}}{ID} &
\multicolumn{1}{>{\raggedright\arraybackslash\color{white}\bfseries}p{2.8cm}}{Nombre} &
\multicolumn{1}{>{\raggedright\arraybackslash\color{white}\bfseries}p{5.2cm}}{Descripción} &
\multicolumn{1}{>{\raggedright\arraybackslash\color{white}\bfseries}p{4.2cm}}{Justificación} \\
\endfirsthead
\rowcolor{azulPUJ}
\multicolumn{1}{>{\centering\arraybackslash\color{white}\bfseries}p{1.5cm}}{ID} &
\multicolumn{1}{>{\raggedright\arraybackslash\color{white}\bfseries}p{2.8cm}}{Nombre} &
\multicolumn{1}{>{\raggedright\arraybackslash\color{white}\bfseries}p{5.2cm}}{Descripción} &
\multicolumn{1}{>{\raggedright\arraybackslash\color{white}\bfseries}p{4.2cm}}{Justificación} \\
\endhead

\rowcolor{grisCl}
\textbf{RD-01} & IEC 61850 &
  Estándar de interoperabilidad para comunicaciones en subestaciones eléctricas mediante
  Ethernet industrial. &
  Garantiza que equipos de distintos fabricantes se comuniquen sin adaptadores propietarios. \\

\textbf{RD-02} & MQTT v5 &
  Protocolo de ingesta de telemetría con mínimo consumo de ancho de banda en redes celulares
  rurales. &
  El protocolo HTTP transmitiría cabeceras pesadas que colapsan conexiones inestables en
  zonas remotas. \\

\rowcolor{grisCl}
\textbf{RD-03} & OPC UA &
  Integración cifrada y segura con los sistemas SCADA legados existentes en las subestaciones. &
  Exponer la base de datos directamente abriría brechas críticas en la red de operación (OT). \\

\textbf{RD-04} & NERC CIP &
  Cumplimiento de la normativa internacional de ciberseguridad para infraestructuras eléctricas
  críticas. &
  Obligación regulatoria; el incumplimiento expone a sanciones legales y riesgo de
  interrupción del servicio nacional. \\

\rowcolor{grisCl}
\textbf{RD-05} & IEEE 1547 &
  Soporte al monitoreo de recursos energéticos distribuidos como paneles solares y baterías. &
  El estándar permite integrar generación renovable sin rediseñar la plataforma. \\

\textbf{RD-06} & PTP IEEE 1588 &
  Sincronización de tiempo con precisión de microsegundos en todos los nodos de la red. &
  Corrige la desincronización de reloj en subestaciones remotas, evitando inconsistencias
  temporales en datos de series de tiempo. \\

\rowcolor{grisCl}
\textbf{RD-07} & ISO 27001 &
  Marco internacional de gestión de seguridad de la información aplicado a todos los controles
  de la plataforma. &
  Proporciona validación auditable de las medidas de seguridad para entes reguladores. \\

\end{longtable}

% ── 1.3 Requerimientos de Producto ───────────────────────────────────────────
\subsection{Requerimientos de Producto}

\begin{longtable}{
  C{1.5cm}
  P{2.8cm}
  P{4.8cm}
  P{4.6cm}
}
\rowcolor{azulPUJ}
\multicolumn{1}{>{\centering\arraybackslash\color{white}\bfseries}p{1.5cm}}{ID} &
\multicolumn{1}{>{\raggedright\arraybackslash\color{white}\bfseries}p{2.8cm}}{Nombre} &
\multicolumn{1}{>{\raggedright\arraybackslash\color{white}\bfseries}p{4.8cm}}{Descripción} &
\multicolumn{1}{>{\raggedright\arraybackslash\color{white}\bfseries}p{4.6cm}}{Criterio de aceptación} \\
\endfirsthead
\rowcolor{azulPUJ}
\multicolumn{1}{>{\centering\arraybackslash\color{white}\bfseries}p{1.5cm}}{ID} &
\multicolumn{1}{>{\raggedright\arraybackslash\color{white}\bfseries}p{2.8cm}}{Nombre} &
\multicolumn{1}{>{\raggedright\arraybackslash\color{white}\bfseries}p{4.8cm}}{Descripción} &
\multicolumn{1}{>{\raggedright\arraybackslash\color{white}\bfseries}p{4.6cm}}{Criterio de aceptación} \\
\endhead

\rowcolor{grisCl}
\textbf{RF-01} & Base de datos de series de tiempo &
  Almacenamiento de millones de lecturas por segundo con timestamp exacto mediante InfluxDB
  Cluster. &
  Escritura sostenida de 50.000 eventos/s sin degradación; consultas por rango en menos de
  2 s. \\

\textbf{RF-02} & Motor de eventos en tiempo real &
  Detección de anomalías de voltaje, vibración, temperatura y presión con disparo de alarmas
  en milisegundos. &
  El 99,9\% de alertas críticas visibles en el panel en menos de 3 s desde el evento. \\

\rowcolor{grisCl}
\textbf{RF-03} & Interfaz geoespacial &
  Panel visual interactivo con ubicación exacta de sensores y fallas sobre mapa nacional en
  tiempo real. &
  El operador localiza cualquier sensor en menos de 5 s de navegación en el mapa. \\

\textbf{RF-04} & Actuación remota con MFA &
  Apagado y reconexión remota de transformadores desde el centro de control con autenticación
  de dos factores. &
  Tiempo de actuación confirmada menor a 10 s; 0 comandos ejecutados sin segundo factor
  válido. \\

\rowcolor{grisCl}
\textbf{RF-05} & Operación offline --- Edge &
  El gateway local procesa alarmas críticas y almacena telemetría durante cortes de
  conectividad a la nube. &
  Cero pérdida de datos tras 24 h de desconexión; resincronización completa en menos de
  15 min. \\

\textbf{RF-06} & Firmware OTA &
  Envío y aplicación de parches de seguridad a los sensores de forma remota sin intervención
  física. &
  100\% de nodos parcheados en menos de 24 h tras publicación de vulnerabilidad crítica. \\

\rowcolor{grisCl}
\textbf{RF-07} & Notificaciones omnicanal &
  Alertas urgentes enviadas por SMS, correo electrónico y notificaciones push a los operadores
  designados. &
  Notificación entregada en menos de 30 s desde la detección en al menos 2 canales
  simultáneos. \\

\textbf{RF-08} & API RESTful de exportación &
  Interfaz para que sistemas de terceros consuman históricos de telemetría con autenticación
  y rate limiting. &
  API disponible al 99,5\% con tiempo de respuesta menor a 3 s para consultas de hasta
  1 millón de registros. \\

\end{longtable}
% ══════════════════════════════════════════════════════════════════════════════
% SECCIÓN 2: USUARIOS
% ══════════════════════════════════════════════════════════════════════════════
\section{Usuarios}

La arquitectura modela diez actores con perfiles diferenciados. El diseño aplica el principio
de menor privilegio sobre cada perfil, garantizando que ningún actor pueda acceder a
funciones fuera de su ámbito operativo. Se distinguen actores humanos y actores máquina
(sistemas externos).

\begin{longtable}{
  P{3.0cm}
  C{1.8cm}
  P{3.2cm}
  P{5.7cm}
}
\rowcolor{azulPUJ}
\multicolumn{1}{>{\raggedright\arraybackslash\color{white}\bfseries}p{3.0cm}}{Actor} &
\multicolumn{1}{>{\centering\arraybackslash\color{white}\bfseries}p{1.8cm}}{Tipo} &
\multicolumn{1}{>{\raggedright\arraybackslash\color{white}\bfseries}p{3.2cm}}{Privilegio} &
\multicolumn{1}{>{\raggedright\arraybackslash\color{white}\bfseries}p{5.7cm}}{Descripción funcional} \\
\endfirsthead
\rowcolor{azulPUJ}
\multicolumn{1}{>{\raggedright\arraybackslash\color{white}\bfseries}p{3.0cm}}{Actor} &
\multicolumn{1}{>{\centering\arraybackslash\color{white}\bfseries}p{1.8cm}}{Tipo} &
\multicolumn{1}{>{\raggedright\arraybackslash\color{white}\bfseries}p{3.2cm}}{Privilegio} &
\multicolumn{1}{>{\raggedright\arraybackslash\color{white}\bfseries}p{5.7cm}}{Descripción funcional} \\
\endhead

\textbf{Operador de Centro de Control} &
  Humano & Visualización + Gestión de alarmas &
  Monitorea tableros en tiempo real y gestiona alarmas operativas. Sin permisos de
  configuración para evitar desconexión accidental de sensores críticos. \\

\rowcolor{grisCl}
\textbf{Ingeniero de Mantenimiento} &
  Humano & Lectura + Registro de campo &
  Acceso móvil al diagnóstico de sensores in situ. Registra calibraciones y cierra
  incidentes físicos en campo. \\

\textbf{Administrador TI/OT} &
  Humano & Configuración + Control de acceso &
  Gestiona la red, emite certificados X.509 y actualiza firmware. Centraliza el control
  de acceso de los gateways industriales. \\

\rowcolor{grisCl}
\textbf{Especialista en Ciberseguridad} &
  Humano & Auditoría + Gestión de políticas &
  Audita tráfico de red, gestiona políticas del WAF/IDS y bloquea nodos comprometidos
  sin afectar la operación eléctrica. \\

\textbf{Analista de Datos Energéticos} &
  Humano & Lectura histórica + Análisis &
  Consulta réplicas de lectura para análisis predictivos pesados sin degradar el
  monitoreo en tiempo real. \\

\rowcolor{grisCl}
\textbf{Gerente de Operaciones} &
  Humano & Lectura ejecutiva &
  Accede a dashboards ejecutivos con métricas de disponibilidad y eficiencia energética.
  No recibe alertas operativas menores. \\

\textbf{Auditor de Cumplimiento} &
  Humano & Solo lectura --- Auditoría &
  Acceso de solo lectura a registros inmutables (WORM) para garantizar transparencia
  ante reguladores sin comprometer datos operativos. \\

\rowcolor{grisCl}
\textbf{Sistema SCADA Legado} &
  Máquina & Recepción de datos consolidados &
  Recibe datos del nuevo sistema vía OPC UA para los paneles de control heredados no
  reemplazables en el corto plazo. \\

\textbf{Sistema Meteorológico Nacional} &
  Máquina & Provisión de datos climáticos &
  Suministra pronósticos vía API REST para cruzar condiciones climáticas con temperatura
  de transformadores y anticipar fallas térmicas. \\

\rowcolor{grisCl}
\textbf{Sistema ERP Corporativo} &
  Máquina & Lectura de métricas de activos &
  Recibe métricas de vida útil de sensores (ej.\ SAP) para automatizar órdenes de
  compra de repuestos. \\

\end{longtable}

% ══════════════════════════════════════════════════════════════════════════════
% SECCIÓN 3: MÓDULOS FUNCIONALES
% ══════════════════════════════════════════════════════════════════════════════
\section{Módulos Funcionales}

La arquitectura divide el sistema en cinco módulos principales alineados con los actores
operativos y con los criterios de calidad de la norma ISO 25010. Cada módulo es desplegable
y escalable de forma independiente sobre Kubernetes.

% ── 3.1 Descripción de Módulos ────────────────────────────────────────────────
\subsection{Descripción de Módulos}

\begin{longtable}{
  P{2.5cm}
  P{4.0cm}
  P{3.5cm}
  P{3.7cm}
}
\rowcolor{azulPUJ}
\multicolumn{1}{>{\raggedright\arraybackslash\color{white}\bfseries}p{2.5cm}}{Módulo} &
\multicolumn{1}{>{\raggedright\arraybackslash\color{white}\bfseries}p{4.0cm}}{Responsabilidad} &
\multicolumn{1}{>{\raggedright\arraybackslash\color{white}\bfseries}p{3.5cm}}{Tecnología principal} &
\multicolumn{1}{>{\raggedright\arraybackslash\color{white}\bfseries}p{3.7cm}}{Dependencias clave} \\
\endfirsthead
\rowcolor{azulPUJ}
\multicolumn{1}{>{\raggedright\arraybackslash\color{white}\bfseries}p{2.5cm}}{Módulo} &
\multicolumn{1}{>{\raggedright\arraybackslash\color{white}\bfseries}p{4.0cm}}{Responsabilidad} &
\multicolumn{1}{>{\raggedright\arraybackslash\color{white}\bfseries}p{3.5cm}}{Tecnología principal} &
\multicolumn{1}{>{\raggedright\arraybackslash\color{white}\bfseries}p{3.7cm}}{Dependencias clave} \\
\endhead

\rowcolor{azulClr}
\textbf{Motor de Alertas} &
  Detecta anomalías en tiempo real mediante reglas CEP y modelos ML. Dispara notificaciones
  omnicanal en menos de 3 s. &
  Apache Flink, Apache Kafka, scikit-learn / TensorFlow Lite &
  Módulo de Ingesta, Base de Datos Hot \\

\rowcolor{azulClr}
\textbf{Panel Geoespacial} &
  Visualiza la ubicación y estado de los 120.000 sensores en mapa nacional interactivo.
  Gestiona alarmas activas por región. &
  Grafana, PostGIS, Leaflet.js, API Gateway &
  Motor de Alertas, Gestión de Nodos \\

\rowcolor{azulClr}
\textbf{Históricos y Analítica} &
  Almacena telemetría a largo plazo con retención de 5 años. Provee capacidades de análisis
  histórico para mantenimiento predictivo. &
  InfluxDB Cluster, Object Storage (S3), Apache Spark &
  Motor de Alertas, API RESTful \\

\rowcolor{azulClr}
\textbf{Gestión de Nodos} &
  Administra el ciclo de vida de los 120.000 sensores y gateways: registro, configuración,
  actualización OTA y monitoreo de salud. &
  Kubernetes Operator, HashiCorp Vault, CI/CD pipeline &
  PKI Interna, Módulo de Auditoría \\

\rowcolor{azulClr}
\textbf{Auditoría y Logs} &
  Registra de forma inmutable todos los accesos, comandos y eventos del sistema para
  cumplimiento regulatorio (NERC CIP, ISO 27001). &
  ELK Stack, WORM Storage, Wazuh SIEM &
  Todos los módulos \\

\end{longtable}

% ── 3.2 Matriz de Trazabilidad: Actores vs. Módulos ──────────────────────────
\subsection{Matriz de Trazabilidad: Actores vs.\ Módulos}

Convenciones: \textbf{C} = Configura, \textbf{U} = Usa, \textbf{L} = Lee,
\textbf{---} = Sin acceso.

\vspace{6pt}

\begin{longtable}{
  P{3.4cm}
  C{1.8cm}
  C{1.8cm}
  C{1.8cm}
  C{1.8cm}
  C{1.8cm}
}
\rowcolor{azulPUJ}
\multicolumn{1}{>{\raggedright\arraybackslash\color{white}\bfseries}p{3.4cm}}{Actor / Módulo} &
\multicolumn{1}{>{\centering\arraybackslash\color{white}\bfseries}p{1.8cm}}{Motor Alertas} &
\multicolumn{1}{>{\centering\arraybackslash\color{white}\bfseries}p{1.8cm}}{Panel Geo.} &
\multicolumn{1}{>{\centering\arraybackslash\color{white}\bfseries}p{1.8cm}}{Históricos} &
\multicolumn{1}{>{\centering\arraybackslash\color{white}\bfseries}p{1.8cm}}{Gestión Nodos} &
\multicolumn{1}{>{\centering\arraybackslash\color{white}\bfseries}p{1.8cm}}{Auditoría} \\
\endfirsthead
\rowcolor{azulPUJ}
\multicolumn{1}{>{\raggedright\arraybackslash\color{white}\bfseries}p{3.4cm}}{Actor / Módulo} &
\multicolumn{1}{>{\centering\arraybackslash\color{white}\bfseries}p{1.8cm}}{Motor Alertas} &
\multicolumn{1}{>{\centering\arraybackslash\color{white}\bfseries}p{1.8cm}}{Panel Geo.} &
\multicolumn{1}{>{\centering\arraybackslash\color{white}\bfseries}p{1.8cm}}{Históricos} &
\multicolumn{1}{>{\centering\arraybackslash\color{white}\bfseries}p{1.8cm}}{Gestión Nodos} &
\multicolumn{1}{>{\centering\arraybackslash\color{white}\bfseries}p{1.8cm}}{Auditoría} \\
\endhead

Operador Centro Control    & \cellcolor{azulClr}U & \cellcolor{azulClr}U & \cellcolor{azulClr}L & ---                  & ---                  \\
\rowcolor{grisCl}
Ing.\ de Mantenimiento     & \cellcolor{azulClr}L & \cellcolor{azulClr}L & \cellcolor{azulClr}L & \cellcolor{azulClr}U & ---                  \\
Administrador TI/OT        & \cellcolor{azulClr}L & \cellcolor{azulClr}L & ---                  & \cellcolor{naranCl}C & \cellcolor{azulClr}L \\
\rowcolor{grisCl}
Esp.\ Ciberseguridad       & ---                  & ---                  & ---                  & \cellcolor{naranCl}C & \cellcolor{azulClr}U \\
Analista de Datos          & ---                  & ---                  & \cellcolor{azulClr}U & ---                  & ---                  \\
\rowcolor{grisCl}
Gerente de Operaciones     & ---                  & \cellcolor{azulClr}L & \cellcolor{azulClr}L & ---                  & ---                  \\
Auditor de Cumplimiento    & ---                  & ---                  & ---                  & ---                  & \cellcolor{azulClr}L \\
\rowcolor{grisCl}
Sistema SCADA Legado       & \cellcolor{azulClr}U & ---                  & ---                  & ---                  & ---                  \\
Sistema ERP Corporativo    & ---                  & ---                  & \cellcolor{azulClr}L & \cellcolor{azulClr}L & ---                  \\
Sistema Meteorológico      & \cellcolor{naranCl}C & ---                  & ---                  & ---                  & ---                  \\

\end{longtable}

% ── 3.3 Matriz de Módulos vs. Atributos ISO 25010 ────────────────────────────
\subsection{Matriz de Módulos vs.\ Atributos ISO 25010}

Niveles de exigencia técnica: \textbf{A} = Alto, \textbf{M} = Medio, \textbf{B} = Bajo.

\vspace{6pt}

\begin{longtable}{
  P{2.8cm}
  C{1.4cm}
  C{1.4cm}
  C{1.4cm}
  C{1.4cm}
  C{1.4cm}
  C{1.4cm}
  C{1.4cm}
}
\rowcolor{azulPUJ}
\multicolumn{1}{>{\raggedright\arraybackslash\color{white}\bfseries}p{2.8cm}}{Módulo} &
\multicolumn{1}{>{\centering\arraybackslash\color{white}\bfseries}p{1.4cm}}{Efic.} &
\multicolumn{1}{>{\centering\arraybackslash\color{white}\bfseries}p{1.4cm}}{Fiabil.} &
\multicolumn{1}{>{\centering\arraybackslash\color{white}\bfseries}p{1.4cm}}{Segur.} &
\multicolumn{1}{>{\centering\arraybackslash\color{white}\bfseries}p{1.4cm}}{Compat.} &
\multicolumn{1}{>{\centering\arraybackslash\color{white}\bfseries}p{1.4cm}}{Usabil.} &
\multicolumn{1}{>{\centering\arraybackslash\color{white}\bfseries}p{1.4cm}}{Mant.} &
\multicolumn{1}{>{\centering\arraybackslash\color{white}\bfseries}p{1.4cm}}{Portab.} \\
\endfirsthead
\rowcolor{azulPUJ}
\multicolumn{1}{>{\raggedright\arraybackslash\color{white}\bfseries}p{2.8cm}}{Módulo} &
\multicolumn{1}{>{\centering\arraybackslash\color{white}\bfseries}p{1.4cm}}{Efic.} &
\multicolumn{1}{>{\centering\arraybackslash\color{white}\bfseries}p{1.4cm}}{Fiabil.} &
\multicolumn{1}{>{\centering\arraybackslash\color{white}\bfseries}p{1.4cm}}{Segur.} &
\multicolumn{1}{>{\centering\arraybackslash\color{white}\bfseries}p{1.4cm}}{Compat.} &
\multicolumn{1}{>{\centering\arraybackslash\color{white}\bfseries}p{1.4cm}}{Usabil.} &
\multicolumn{1}{>{\centering\arraybackslash\color{white}\bfseries}p{1.4cm}}{Mant.} &
\multicolumn{1}{>{\centering\arraybackslash\color{white}\bfseries}p{1.4cm}}{Portab.} \\
\endhead

Motor de Alertas  & \cellcolor{naranCl}A & \cellcolor{naranCl}A & \cellcolor{naranCl}A & \cellcolor{grisCl}M  & B                    & \cellcolor{grisCl}M  & \cellcolor{grisCl}M  \\
\rowcolor{grisCl}
Panel Geoespacial & \cellcolor{naranCl}A & \cellcolor{grisCl}M  & \cellcolor{grisCl}M  & \cellcolor{grisCl}M  & \cellcolor{naranCl}A & \cellcolor{grisCl}M  & \cellcolor{naranCl}A \\
Históricos        & \cellcolor{naranCl}A & \cellcolor{naranCl}A & \cellcolor{grisCl}M  & B                    & B                    & \cellcolor{grisCl}M  & B                    \\
\rowcolor{grisCl}
Gestión de Nodos  & \cellcolor{grisCl}M  & \cellcolor{naranCl}A & \cellcolor{naranCl}A & \cellcolor{naranCl}A & \cellcolor{grisCl}M  & \cellcolor{naranCl}A & \cellcolor{naranCl}A \\
Auditoría y Logs  & \cellcolor{grisCl}M  & \cellcolor{naranCl}A & \cellcolor{naranCl}A & \cellcolor{grisCl}M  & B                    & \cellcolor{naranCl}A & B                    \\

\end{longtable}
% ══════════════════════════════════════════════════════════════════════════════
% SECCIÓN 4: RESTRICCIONES Y RIESGOS
% ══════════════════════════════════════════════════════════════════════════════
\section{Restricciones y Riesgos}

El diseño arquitectónico asume limitaciones innegociables impuestas por el entorno operativo
y anticipa problemas futuros que amenazan la viabilidad del sistema.

% ── 4.1 Restricciones Arquitectónicas ────────────────────────────────────────
\subsection{Restricciones Arquitectónicas}

\begin{longtable}{
  C{1.5cm}
  P{2.8cm}
  P{4.8cm}
  P{4.6cm}
}
\rowcolor{azulPUJ}
\multicolumn{1}{>{\centering\arraybackslash\color{white}\bfseries}p{1.5cm}}{ID} &
\multicolumn{1}{>{\raggedright\arraybackslash\color{white}\bfseries}p{2.8cm}}{Restricción} &
\multicolumn{1}{>{\raggedright\arraybackslash\color{white}\bfseries}p{4.8cm}}{Descripción} &
\multicolumn{1}{>{\raggedright\arraybackslash\color{white}\bfseries}p{4.6cm}}{Impacto arquitectónico} \\
\endfirsthead
\rowcolor{azulPUJ}
\multicolumn{1}{>{\centering\arraybackslash\color{white}\bfseries}p{1.5cm}}{ID} &
\multicolumn{1}{>{\raggedright\arraybackslash\color{white}\bfseries}p{2.8cm}}{Restricción} &
\multicolumn{1}{>{\raggedright\arraybackslash\color{white}\bfseries}p{4.8cm}}{Descripción} &
\multicolumn{1}{>{\raggedright\arraybackslash\color{white}\bfseries}p{4.6cm}}{Impacto arquitectónico} \\
\endhead

\rowcolor{grisCl}
\textbf{RC-01} & Integración obligada con SCADA legado &
  El cliente prohíbe reemplazar los controladores físicos actuales. La nueva plataforma
  debe conectarse mediante OPC UA. &
  Obliga a implementar una capa de adaptación OPC UA y limita la frecuencia de
  actualización de datos heredados. \\

\textbf{RC-02} & Presupuesto limitado para licencias &
  El proyecto carece de fondos para licencias corporativas por sensor. Toda la pila
  tecnológica debe ser open-source. &
  Tecnologías como Oracle DB, Splunk o Azure Event Hubs quedan descartadas; la elección
  recae en InfluxDB, Kafka y Grafana. \\

\rowcolor{grisCl}
\textbf{RC-03} & Hardware de borde con recursos limitados &
  Las subestaciones antiguas no tienen espacio ni refrigeración para servidores de
  tamaño completo. Los gateways deben operar en hardware ARM de baja potencia. &
  Los algoritmos de detección local deben ser ligeros; modelos ML complejos ejecutan
  en la nube, no en el borde. \\

\textbf{RC-04} & Retención regulatoria de 5 años &
  La normativa exige conservar datos operativos históricos por al menos 5 años con
  integridad garantizable. &
  Requiere estrategia de almacenamiento por niveles (hot/warm/cold) y uso de WORM
  storage para inmutabilidad. \\

\rowcolor{grisCl}
\textbf{RC-05} & Disponibilidad 99,95\% anual &
  El SLA exige un máximo de 4,38 h de inactividad acumulada por año, incluyendo
  mantenimientos programados. &
  Impone arquitectura multi-zona activa-activa con failover automático y despliegues
  blue/green sin downtime. \\

\end{longtable}

% ── 4.2 Riesgos del Proyecto ─────────────────────────────────────────────────
\subsection{Riesgos del Proyecto}

\begin{longtable}{
  C{1.2cm}
  P{2.8cm}
  C{1.2cm}
  C{1.4cm}
  P{3.6cm}
  P{3.5cm}
}
\rowcolor{azulPUJ}
\multicolumn{1}{>{\centering\arraybackslash\color{white}\bfseries}p{1.2cm}}{ID} &
\multicolumn{1}{>{\raggedright\arraybackslash\color{white}\bfseries}p{2.8cm}}{Riesgo} &
\multicolumn{1}{>{\centering\arraybackslash\color{white}\bfseries}p{1.2cm}}{Prob.} &
\multicolumn{1}{>{\centering\arraybackslash\color{white}\bfseries}p{1.4cm}}{Impacto} &
\multicolumn{1}{>{\raggedright\arraybackslash\color{white}\bfseries}p{3.6cm}}{Plan de mitigación} &
\multicolumn{1}{>{\raggedright\arraybackslash\color{white}\bfseries}p{3.5cm}}{Contingencia} \\
\endfirsthead
\rowcolor{azulPUJ}
\multicolumn{1}{>{\centering\arraybackslash\color{white}\bfseries}p{1.2cm}}{ID} &
\multicolumn{1}{>{\raggedright\arraybackslash\color{white}\bfseries}p{2.8cm}}{Riesgo} &
\multicolumn{1}{>{\centering\arraybackslash\color{white}\bfseries}p{1.2cm}}{Prob.} &
\multicolumn{1}{>{\centering\arraybackslash\color{white}\bfseries}p{1.4cm}}{Impacto} &
\multicolumn{1}{>{\raggedright\arraybackslash\color{white}\bfseries}p{3.6cm}}{Plan de mitigación} &
\multicolumn{1}{>{\raggedright\arraybackslash\color{white}\bfseries}p{3.5cm}}{Contingencia} \\
\endhead

\textbf{RZ-01} &
  Saturación en capa de ingesta por picos extremos (+60\% volumen) &
  \cellcolor{naranCl}Alta & \cellcolor{red!15}Crítico &
  Auto-escalado horizontal en Kafka con HPA de Kubernetes basado en lag del consumer
  group. &
  Circuit breaker en gateways edge para priorizar eventos críticos y descartar
  telemetría no urgente. \\

\rowcolor{grisCl}
\textbf{RZ-02} &
  Pérdida de conectividad en nodos remotos durante eventos climáticos &
  \cellcolor{naranCl}Alta & \cellcolor{naranCl}Alto &
  Edge Computing con almacenamiento local de 24 h y mecanismo Store-and-Forward
  con MQTT QoS~2. &
  Activación de enlace 4G/LTE de respaldo; alertas locales sin dependencia de nube
  central. \\

\textbf{RZ-03} &
  Inconsistencias temporales por desincronización de reloj en regiones &
  \cellcolor{verdeCl}Media & \cellcolor{naranCl}Alto &
  Protocolo PTP IEEE~1588 con Grandmaster Clock por región; corrección de clock skew
  en Flink. &
  Marcas de tiempo relativas al gateway local + reconciliación post-hoc con NTP como
  fallback. \\

\rowcolor{grisCl}
\textbf{RZ-04} &
  Ciberataque a infraestructura crítica (inyección de datos o comandos) &
  \cellcolor{verdeCl}Media & \cellcolor{red!15}Crítico &
  mTLS con X.509 por dispositivo, segmentación de red OT/IT, WAF, IDS/IPS y
  auditoría en tiempo real. &
  Aislamiento inmediato del nodo comprometido; activación de modo degradado con
  operación manual local. \\

\textbf{RZ-05} &
  Cuellos de botella en procesamiento streaming durante tormentas &
  \cellcolor{verdeCl}Media & \cellcolor{naranCl}Alto &
  Particionado de Kafka por región geográfica; Flink con backpressure automático y
  ventanas de tiempo adaptativas. &
  Reducción dinámica de granularidad de sensores no críticos; procesamiento diferido
  para análisis no urgente. \\

\rowcolor{grisCl}
\textbf{RZ-06} &
  Costos exponenciales de almacenamiento a largo plazo &
  \cellcolor{verdeCl}Media & \cellcolor{grisCl}Medio &
  Política de retención en InfluxDB (90 días hot), compresión warm tier, Object
  Storage para cold ($>$1 año). &
  Revisión anual de política de retención; downsampling de métricas de baja prioridad
  después de 6 meses. \\

\textbf{RZ-07} &
  Vulneración física de gateways en subestaciones sin vigilancia &
  \cellcolor{verdeCl}Baja & \cellcolor{red!15}Crítico &
  Cifrado de disco AES-256, TPM~2.0, detección de apertura de carcasa (tamper
  detection), VPN site-to-site. &
  Revocación remota de certificados X.509 del dispositivo comprometido; reemplazo de
  hardware programado. \\

\end{longtable}
% ══════════════════════════════════════════════════════════════════════════════
% SECCIÓN 5: ATRIBUTOS DE CALIDAD SEGÚN ISO 25010
% ══════════════════════════════════════════════════════════════════════════════
\section{Atributos de Calidad según ISO 25010}

La norma ISO 25010 define ocho características de calidad del producto software. La
arquitectura prioriza \textbf{Eficiencia de Desempeño}, \textbf{Fiabilidad} y
\textbf{Seguridad} como atributos críticos (Nivel~1) dado el contexto de infraestructura
energética nacional. Los cinco restantes se clasifican como importantes (Nivel~2) y están
subordinados a los pilares críticos.

\begin{longtable}{
  C{1.2cm}
  P{2.8cm}
  C{1.8cm}
  P{3.8cm}
  P{4.1cm}
}
\rowcolor{azulPUJ}
\multicolumn{1}{>{\centering\arraybackslash\color{white}\bfseries}p{1.2cm}}{Nivel} &
\multicolumn{1}{>{\raggedright\arraybackslash\color{white}\bfseries}p{2.8cm}}{Atributo} &
\multicolumn{1}{>{\centering\arraybackslash\color{white}\bfseries}p{1.8cm}}{Prioridad} &
\multicolumn{1}{>{\raggedright\arraybackslash\color{white}\bfseries}p{3.8cm}}{Razón de priorización} &
\multicolumn{1}{>{\raggedright\arraybackslash\color{white}\bfseries}p{4.1cm}}{Impacto operativo si se incumple} \\
\endfirsthead
\rowcolor{azulPUJ}
\multicolumn{1}{>{\centering\arraybackslash\color{white}\bfseries}p{1.2cm}}{Nivel} &
\multicolumn{1}{>{\raggedright\arraybackslash\color{white}\bfseries}p{2.8cm}}{Atributo} &
\multicolumn{1}{>{\centering\arraybackslash\color{white}\bfseries}p{1.8cm}}{Prioridad} &
\multicolumn{1}{>{\raggedright\arraybackslash\color{white}\bfseries}p{3.8cm}}{Razón de priorización} &
\multicolumn{1}{>{\raggedright\arraybackslash\color{white}\bfseries}p{4.1cm}}{Impacto operativo si se incumple} \\
\endhead

\cellcolor{naranCl}\textbf{1} & \textbf{Eficiencia de Desempeño} &
  \cellcolor{red!15}Crítica &
  El sistema procesa 2,4~M de eventos/min con SLA de latencia de 3~s para alertas
  críticas. Cualquier degradación es inaceptable en tiempo real. &
  Una latencia mayor a 3~s impide detectar sobrecargas en transformadores a tiempo,
  causando explosiones físicas y apagones regionales. \\

\rowcolor{grisCl}
\cellcolor{naranCl}\textbf{1} & \textbf{Fiabilidad} &
  \cellcolor{red!15}Crítica &
  La red eléctrica exige monitoreo continuo 24/7/365. Un tiempo de inactividad superior
  a 4,38~h/año viola el SLA contractual. &
  La caída del sistema deja ciegos a los centros de control y retrasa la respuesta ante
  apagones en cascada a nivel ciudad o regional. \\

\cellcolor{naranCl}\textbf{1} & \textbf{Seguridad} &
  \cellcolor{red!15}Crítica &
  La ley NERC~CIP obliga a proteger infraestructuras eléctricas. La plataforma es
  objetivo directo de actores estatales y ciberterroristas. &
  Una vulnerabilidad exitosa permite apagar subestaciones remotamente, paralizando
  hospitales, acueductos y sistemas de comunicación. \\

\rowcolor{grisCl}
\cellcolor{verdeCl}\textbf{2} & \textbf{Compatibilidad} &
  \cellcolor{verdeCl}Importante &
  Coexistencia obligada con sistemas SCADA legados y estándares industriales
  IEC~61850 y OPC~UA. &
  Sin compatibilidad, el sistema operaría en silos sin visibilidad sobre equipos
  heredados que representan el 80\% de la infraestructura actual. \\

\cellcolor{verdeCl}\textbf{2} & \textbf{Usabilidad} &
  \cellcolor{verdeCl}Importante &
  Operadores de centro de control deben actuar en segundos bajo presión. Errores de
  interfaz son inaceptables en emergencias. &
  Una interfaz confusa puede llevar a un operador a desconectar el sensor equivocado o
  ignorar una alarma crítica por fatiga visual. \\

\rowcolor{grisCl}
\cellcolor{verdeCl}\textbf{2} & \textbf{Mantenibilidad} &
  \cellcolor{verdeCl}Importante &
  La arquitectura de microservicios permite actualizar módulos sin detener el
  monitoreo. Esencial para parches de seguridad urgentes. &
  Sin separación de módulos, cada actualización requiere ventana de mantenimiento
  total, violando el SLA de disponibilidad. \\

\cellcolor{verdeCl}\textbf{2} & \textbf{Portabilidad} &
  \cellcolor{verdeCl}Importante &
  El empaquetamiento en Docker/Kubernetes permite migrar entre proveedores de nube
  para evitar vendor lock-in regulatorio. &
  La dependencia de un único proveedor cloud expone la infraestructura crítica a
  fallos de disponibilidad del proveedor. \\

\rowcolor{grisCl}
\cellcolor{verdeCl}\textbf{2} & \textbf{Adecuación Funcional} &
  \cellcolor{verdeCl}Importante &
  El sistema debe capturar métricas eléctricas precisas y generar alertas correctas
  para los tipos de sensores del dominio energético. &
  Funciones incompletas generan falsos negativos (anomalías no detectadas) o falsos
  positivos que saturan a los operadores. \\

\end{longtable}
% ══════════════════════════════════════════════════════════════════════════════
% SECCIÓN 6: ESCENARIOS DE CALIDAD
% ══════════════════════════════════════════════════════════════════════════════
\section{Escenarios de Calidad}

Se definen escenarios de calidad para las subcategorías más críticas de ISO~25010,
utilizando la plantilla estándar de seis elementos: \textbf{Fuente}, \textbf{Estímulo},
\textbf{Artefacto}, \textbf{Entorno}, \textbf{Respuesta} y \textbf{Medida de respuesta}.

% Macro para uniformizar la tabla de cada escenario
\newcommand{\escenario}[7]{%
  \begin{longtable}{P{2.8cm} P{10.9cm}}
  \rowcolor{azulPUJ}
  \multicolumn{2}{l}{%
    \textbf{\color{white}#1}%
  } \\
  \endfirsthead
  \rowcolor{azulPUJ}
  \multicolumn{2}{l}{%
    \textbf{\color{white}#1}%
  } \\
  \endhead
  \rowcolor{azulClr}\textbf{Fuente}    & #2 \\
  \textbf{Estímulo}                    & #3 \\
  \rowcolor{azulClr}\textbf{Artefacto} & #4 \\
  \textbf{Entorno}                     & #5 \\
  \rowcolor{azulClr}\textbf{Respuesta} & #6 \\
  \textbf{Medida}                      & #7 \\
  \end{longtable}
  \vspace{4pt}
}

% ── 6.1 Eficiencia de Desempeño ───────────────────────────────────────────────
\subsection{Eficiencia de Desempeño}

\escenario{ESC-01 \textnormal{---} Comportamiento temporal bajo carga extrema}
  {Red completa de 120.000 sensores IoT desplegados en subestaciones nacionales.}
  {Los sensores envían telemetría simultánea debido a una tormenta eléctrica que
   incrementa el volumen en un 60\% (hasta 3,84~M eventos/min).}
  {Clúster de ingesta Kafka + Motor de procesamiento Apache Flink.}
  {Operación degradada por evento climático extremo; hora pico de consumo eléctrico.}
  {El sistema activa auto-escalado horizontal en Kafka (nuevas particiones) y en Flink
   (nuevos task managers). Los eventos se encolan sin rechazo de conexiones. El motor
   prioriza eventos críticos sobre telemetría de baja prioridad mediante backpressure
   adaptativo.}
  {El 99\% de las alertas críticas aparecen en el panel del operador en menos de 3~s
   desde la generación en el sensor. Throughput sostenido de al menos 3,84~M
   eventos/min sin degradación de latencia.}

\escenario{ESC-02 \textnormal{---} Utilización de recursos en gateways edge}
  {Edge Gateway ARM con 4~GB de RAM operando en subestación remota.}
  {1.000 sensores locales envían datos simultáneamente durante un pico de voltaje zonal.}
  {Motor de reglas local en el Edge Gateway (InfluxDB local + reglas CEP simplificadas).}
  {Operación normal; conectividad con la nube disponible; temperatura ambiente dentro
   de rango operativo.}
  {El gateway evalúa las reglas críticas localmente, almacena telemetría en buffer local
   y reenvía a la nube sin superar el 70\% de CPU ni el 80\% de RAM.}
  {CPU promedio menor al 70\%; RAM menor al 80\%; tiempo de evaluación de reglas
   críticas menor a 500~ms en el dispositivo.}

% ── 6.2 Fiabilidad ────────────────────────────────────────────────────────────
\subsection{Fiabilidad}

\escenario{ESC-03 \textnormal{---} Tolerancia a fallos --- Corte de conectividad regional}
  {Fenómeno climático severo (huracán o tormenta eléctrica en la Costa Caribe
   colombiana).}
  {Se corta el enlace principal de fibra óptica entre 15 subestaciones rurales y el
   centro de datos regional. El enlace 4G de respaldo también cae por saturación de la
   red móvil.}
  {Edge Gateways instalados en cada subestación afectada.}
  {Operación degradada --- desconexión total de la nube central por más de 24 horas.}
  {Los gateways continúan evaluando reglas críticas de forma autónoma, ejecutando
   apagados de emergencia locales si detectan sobrecarga, y almacenando telemetría en
   InfluxDB local. Al restaurarse la conexión, sincronizan datos mediante MQTT QoS~2
   con reintentos automáticos.}
  {Cero pérdida de datos críticos tras 24~h de desconexión. Sincronización completa
   con la nube en menos de 15~min tras restaurar el enlace. El operador no pierde
   trazabilidad histórica de ningún evento ocurrido durante la desconexión.}

\escenario{ESC-04 \textnormal{---} Recuperabilidad --- Fallo catastrófico del data center}
  {Fallo de infraestructura eléctrica en el data center principal (zona de
   disponibilidad primaria).}
  {Caída total de todos los servicios cloud en la zona primaria, incluyendo bases de
   datos, procesamiento y API gateway.}
  {Arquitectura multi-zona activa-activa con replicación cross-region continua.}
  {Desastre total de la zona primaria; zona secundaria en estado warm standby.}
  {El balanceador de carga global redirige automáticamente el tráfico a la zona
   secundaria. Kubernetes detecta la indisponibilidad y activa los pods en la zona
   alternativa. Los datos replicados en tiempo real garantizan consistencia.}
  {RTO menor a 15~min. RPO menor a 30~s de datos perdidos. Disponibilidad anual
   resultante mayor al 99,95\%.}

% ── 6.3 Seguridad ─────────────────────────────────────────────────────────────
\subsection{Seguridad}

\escenario{ESC-05 \textnormal{---} Confidencialidad --- Intercepción de tráfico OT}
  {Atacante externo con acceso a la red de telecomunicaciones que transporta tráfico
   de la subestación.}
  {El atacante realiza un ataque man-in-the-middle para interceptar y decodificar
   telemetría de la red eléctrica.}
  {Canal de comunicación MQTT entre sensores/gateways y el MQTT Broker en la nube.}
  {Operación normal; la red de telecomunicaciones está comprometida en un tramo
   intermedio.}
  {Todo el tráfico MQTT viaja cifrado con TLS~1.3. La autenticación mutua mTLS con
   certificados X.509 únicos por dispositivo impide que el atacante se identifique como
   un nodo legítimo. El tráfico interceptado es ilegible sin las claves privadas del
   dispositivo, almacenadas en el TPM~2.0.}
  {0\% de tráfico sensible accesible en texto plano. Tiempo de detección de dispositivo
   con certificado inválido menor a 50~ms. El intento de conexión queda registrado en
   el SIEM.}

\escenario{ESC-06 \textnormal{---} Integridad --- Inyección de comandos maliciosos}
  {Atacante que ha comprometido credenciales de un operador mediante phishing.}
  {El atacante intenta enviar comandos de apagado a transformadores de alta tensión
   desde un dispositivo no autorizado.}
  {Módulo de Gestión de Nodos --- API Gateway de actuación remota.}
  {Operación bajo ataque activo; credenciales del operador comprometidas.}
  {El sistema exige autenticación multifactor (MFA) para comandos de actuación. El
   certificado del dispositivo del atacante no coincide con el registrado. El WAF
   detecta el patrón de acceso inusual (horario, geolocalización) y bloquea la sesión.
   El SIEM genera alerta inmediata al Especialista en Ciberseguridad.}
  {100\% de comandos de actuación requieren MFA válido. El sistema bloquea la sesión
   comprometida en menos de 5~s tras detección de anomalía. El intento queda registrado
   con hash inmutable en el log de auditoría.}

% ── 6.4 Compatibilidad ────────────────────────────────────────────────────────
\subsection{Compatibilidad}

\escenario{ESC-07 \textnormal{---} Interoperabilidad con sistemas SCADA legados}
  {Centro de control regional con sistema SCADA de 15 años de antigüedad.}
  {El SCADA legado solicita datos consolidados de voltaje y temperatura de las
   subestaciones de su región en el formato propietario que entiende.}
  {Capa de adaptación OPC~UA del módulo de Gestión de Nodos.}
  {Operación normal; el sistema legado no ha sido modificado ni actualizado.}
  {El adaptador OPC~UA traduce los datos del formato interno de la plataforma al modelo
   de información IEC~61850 que el SCADA legado consume. La comunicación es
   bidireccional y cifrada.}
  {100\% de los sistemas SCADA registrados reciben datos en su formato esperado.
   Latencia de traducción OPC~UA menor a 1~s. Sin modificaciones al software legado.}

% ── 6.5 Usabilidad ────────────────────────────────────────────────────────────
\subsection{Usabilidad}

\escenario{ESC-08 \textnormal{---} Prevención de errores operativos bajo presión}
  {Operador de Centro de Control durante un blackout regional con 200 alarmas activas
   simultáneas.}
  {El operador debe identificar el sensor con la anomalía más crítica y ejecutar un
   apagado de emergencia en menos de 2~minutos.}
  {Panel Geoespacial --- interfaz de gestión de alarmas.}
  {Emergencia operativa; múltiples alarmas simultáneas; presión alta sobre el operador.}
  {El panel ordena automáticamente las alarmas por criticidad y geolocalización. El
   mapa resalta visualmente el transformador en riesgo. Un diálogo de confirmación con
   countdown previene actuaciones accidentales.}
  {El operador localiza el sensor crítico en menos de 30~s. Tasa de actuaciones
   accidentales (comandos sobre el nodo equivocado) menor al 0,1\%.}

% ── 6.6 Mantenibilidad ────────────────────────────────────────────────────────
\subsection{Mantenibilidad}

\escenario{ESC-09 \textnormal{---} Modificabilidad --- Actualización de módulo sin downtime}
  {Equipo de desarrollo que detecta una vulnerabilidad crítica en el Motor de Alertas.}
  {Se requiere desplegar un parche de seguridad en producción dentro de las 4~horas
   siguientes a la publicación del CVE.}
  {Motor de Alertas --- microservicio independiente en Kubernetes.}
  {Sistema en producción con monitoreo activo de 120.000 sensores.}
  {El pipeline CI/CD ejecuta pruebas automáticas del parche en staging en menos de
   20~min. Kubernetes realiza un rolling update del Motor de Alertas sin afectar los
   demás módulos. El tráfico se redirige gradualmente a las nuevas réplicas.}
  {Tiempo total del parche (commit $\rightarrow$ producción) menor a 2~h. Cero
   interrupciones del monitoreo durante el despliegue. 100\% de las pruebas de
   regresión automáticas pasadas antes del despliegue.}

% ── 6.7 Portabilidad ──────────────────────────────────────────────────────────
\subsection{Portabilidad}

\escenario{ESC-10 \textnormal{---} Adaptabilidad --- Migración entre proveedores cloud}
  {Decisión regulatoria que exige mover la infraestructura a un proveedor cloud con
   sede en Colombia.}
  {El equipo de infraestructura recibe la orden de migrar la plataforma completa a un
   nuevo proveedor cloud en un plazo de 6~meses.}
  {Plataforma completa empaquetada en contenedores Docker orquestada con Kubernetes.}
  {Operación continua requerida durante la migración; sin ventana de mantenimiento
   total disponible.}
  {La arquitectura cloud-agnostic con Kubernetes, Helm charts y Terraform permite
   reprovisionar el entorno en el nuevo proveedor con mínimas modificaciones. La
   migración de datos usa herramientas estándar de InfluxDB y PostgreSQL.}
  {Migración completa en menos de 6~meses. Sin reescritura de lógica de negocio. Menos
   de 4~h de downtime total durante el proceso de migración.}

% ── 6.8 Adecuación Funcional ──────────────────────────────────────────────────
\subsection{Adecuación Funcional}

\escenario{ESC-11 \textnormal{---} Corrección funcional --- Precisión en detección de anomalías}
  {Sensor de voltaje en subestación de alta tensión.}
  {El sensor reporta una lectura de voltaje un 15\% por encima del umbral nominal
   (indicador temprano de sobrecarga antes del fallo físico).}
  {Motor de Alertas --- reglas CEP y modelo de scoring ML.}
  {Operación normal; el evento ocurre en horario de baja carga administrativa.}
  {El Motor de Alertas evalúa la lectura contra las reglas de umbral configuradas y el
   modelo de predicción de fallo. Clasifica el evento como Alerta Crítica Nivel~1 y
   la envía al Operador del Centro de Control y al Ingeniero de Mantenimiento
   responsable de la zona.}
  {Tasa de detección verdadera (recall) mayor al 99,5\% para anomalías de nivel
   crítico. Tasa de falsos positivos menor al 2\%. Clasificación correcta del nivel de
   criticidad en más del 98\% de los eventos.}

   % ══════════════════════════════════════════════════════════════════════════════
% SECCIÓN 7: ACUERDOS DE LOS ATRIBUTOS DE CALIDAD (SLA)
% ══════════════════════════════════════════════════════════════════════════════
\section{Acuerdos de los Atributos de Calidad (SLA)}

Los Acuerdos de Nivel de Servicio formalizan las métricas de calidad que el sistema debe
cumplir de forma continua y verificable. El equipo técnico monitorea estas métricas mediante
dashboards de Prometheus y Grafana, generando alertas automáticas ante desviaciones.

\begin{longtable}{
  C{1.2cm}
  P{2.2cm}
  P{2.2cm}
  P{4.8cm}
  P{3.3cm}
}
\rowcolor{azulPUJ}
\multicolumn{1}{>{\centering\arraybackslash\color{white}\bfseries}p{1.2cm}}{ID} &
\multicolumn{1}{>{\raggedright\arraybackslash\color{white}\bfseries}p{2.2cm}}{Atributo} &
\multicolumn{1}{>{\raggedright\arraybackslash\color{white}\bfseries}p{2.2cm}}{Métrica} &
\multicolumn{1}{>{\raggedright\arraybackslash\color{white}\bfseries}p{4.8cm}}{Acuerdo (SLA)} &
\multicolumn{1}{>{\raggedright\arraybackslash\color{white}\bfseries}p{3.3cm}}{Herramienta de medición} \\
\endfirsthead
\rowcolor{azulPUJ}
\multicolumn{1}{>{\centering\arraybackslash\color{white}\bfseries}p{1.2cm}}{ID} &
\multicolumn{1}{>{\raggedright\arraybackslash\color{white}\bfseries}p{2.2cm}}{Atributo} &
\multicolumn{1}{>{\raggedright\arraybackslash\color{white}\bfseries}p{2.2cm}}{Métrica} &
\multicolumn{1}{>{\raggedright\arraybackslash\color{white}\bfseries}p{4.8cm}}{Acuerdo (SLA)} &
\multicolumn{1}{>{\raggedright\arraybackslash\color{white}\bfseries}p{3.3cm}}{Herramienta de medición} \\
\endhead

\rowcolor{grisCl}
\textbf{SLA-01} & Eficiencia de Desempeño & Latencia de alertas críticas &
  El 99,9\% de las alertas críticas aparecen en el panel del operador en menos de 3~s
  desde la detección física en el sensor, medido de extremo a extremo incluyendo
  transporte y procesamiento. &
  Prometheus + Grafana con histogramas de latencia percentil p99. \\

\textbf{SLA-02} & Eficiencia de Desempeño & Throughput de ingesta &
  El sistema ingesta y procesa un mínimo de 50.000 eventos por segundo de forma
  sostenida. Durante picos extremos (+60\%), el throughput escala a 80.000
  eventos/s sin rechazar conexiones. &
  Kafka Consumer Lag Monitor + Flink Dashboard. \\

\rowcolor{grisCl}
\textbf{SLA-03} & Fiabilidad & Disponibilidad (Uptime) &
  La plataforma central mantiene una disponibilidad operativa del 99,95\% anual,
  equivalente a un máximo de 4,38~horas de inactividad acumulada por año, incluyendo
  mantenimientos programados. &
  Prometheus Alertmanager + Uptime Robot externo. \\

\textbf{SLA-04} & Fiabilidad & Tiempo de recuperación (RTO) &
  El sistema restaura la operación completa en la zona secundaria en menos de
  15~minutos tras una falla catastrófica del centro de datos principal. El RPO
  (pérdida máxima de datos) es de 30~segundos. &
  Pruebas de failover trimestrales + métricas de replicación cross-region. \\

\rowcolor{grisCl}
\textbf{SLA-05} & Fiabilidad & Integridad de datos offline &
  Cero pérdida de eventos críticos tras un corte de conectividad de hasta 24~horas
  en nodos edge. Los datos se sincronizan completamente en menos de 15~minutos al
  restaurar el enlace. &
  Mecanismo de reconciliación MQTT QoS~2 + alertas de buffer overflow en gateway. \\

\textbf{SLA-06} & Seguridad & Bloqueo de intrusiones &
  El sistema bloquea el 100\% de las solicitudes con certificado digital inválido o
  ausente en menos de 50~ms. El WAF bloquea el 99,9\% de ataques conocidos
  (OWASP Top~10 industrial). &
  Wazuh SIEM + dashboards de cortafuegos perimetral. \\

\rowcolor{grisCl}
\textbf{SLA-07} & Seguridad & Tiempo de parcheo &
  El equipo aplica parches críticos (CVSS $\geq$ 9.0) en todos los nodos en un plazo
  máximo de 24~horas tras la publicación del CVE. Para severidad alta
  (CVSS 7.0--8.9), el plazo es 72~horas. &
  Pipeline CI/CD con OTA Manager + registro de versiones por dispositivo. \\

\textbf{SLA-08} & Compatibilidad & Disponibilidad OPC~UA &
  La capa de integración con SCADA legados tiene disponibilidad del 99,9\% anual.
  Los datos se entregan a los sistemas heredados con latencia menor a 1~s. &
  Monitoring de endpoints OPC~UA + alertas de desconexión de SCADA registrado. \\

\rowcolor{grisCl}
\textbf{SLA-09} & Mantenibilidad & Lead Time de despliegue &
  El tiempo desde que un commit llega a la rama principal hasta que está en producción
  es menor a 2~horas para parches críticos y menor a 24~horas para funcionalidades
  nuevas. &
  ArgoCD pipeline metrics + DORA dashboard en Grafana. \\

\textbf{SLA-10} & Almacenamiento & Retención de históricos &
  Los datos operativos se conservan durante un mínimo de 5~años con integridad
  verificable. El crecimiento proyectado de 80~TB/año se gestiona automáticamente
  mediante la política de niveles hot/warm/cold. &
  InfluxDB retention policies + alertas de capacidad de Object Storage. \\

\end{longtable}

% ══════════════════════════════════════════════════════════════════════════════
% SECCIÓN 8: REGISTROS DE DECISIONES ARQUITECTÓNICAS (ADRs)
% ══════════════════════════════════════════════════════════════════════════════
\section{Registros de Decisiones Arquitectónicas (ADRs)}

Los ADRs documentan las decisiones estructurales más importantes de la arquitectura,
incluyendo el contexto que las motivó, las alternativas consideradas y descartadas,
la decisión tomada, y las consecuencias positivas y negativas resultantes.

% Macro para uniformizar cada ADR
\newcommand{\adr}[8]{%
  \begin{longtable}{P{3.2cm} P{10.5cm}}
  \rowcolor{azulPUJ}
  \multicolumn{2}{l}{\textbf{\color{white}#1}} \\
  \endfirsthead
  \rowcolor{azulPUJ}
  \multicolumn{2}{l}{\textbf{\color{white}#1}} \\
  \endhead
  \rowcolor{azulClr}  \textbf{Estado}                    & #2 \\
                      \textbf{Contexto}                   & #3 \\
  \rowcolor{azulClr}  \textbf{Alternativas consideradas}  & #4 \\
  \rowcolor{naranCl}  \textbf{Decisión}                   & #5 \\
  \rowcolor{verdeCl}  \textbf{Consecuencias (+)}          & #6 \\
  \rowcolor{red!8}    \textbf{Consecuencias (--)}         & #7 \\
  \rowcolor{azulClr}  \textbf{Cómo validar}               & #8 \\
  \end{longtable}
  \vspace{4pt}
}

% ── ADR-001 ───────────────────────────────────────────────────────────────────
\adr{ADR-001 \textnormal{---} Arquitectura de Microservicios sobre Kubernetes}
  {Aprobado}
  {El sistema procesa 50.000 eventos/s de forma sostenida y sirve paneles en tiempo
   real a múltiples usuarios simultáneamente. Un fallo en el módulo de alertas no debe
   afectar la visualización del mapa geoespacial ni el módulo histórico.}
  {(1)~\textbf{Monolito modular}: simple de desplegar, pero un fallo en cualquier
   módulo tumba el sistema completo --- inaceptable para el SLA de 99,95\%.
   (2)~\textbf{SOA con ESB}: mayor overhead de comunicación y punto único de fallo en
   el bus.
   (3)~\textbf{Microservicios con Kubernetes}: aislamiento de fallos, escalado
   independiente por módulo y gestión declarativa del estado.}
  {Se adopta arquitectura de microservicios empaquetados en contenedores Docker y
   orquestados con Kubernetes. Cada uno de los cinco módulos funcionales es un servicio
   independiente con su propio ciclo de despliegue.}
  {Escalado independiente del Motor de Alertas sin afectar la interfaz web. Despliegues
   rolling update sin downtime. Aislamiento de fallos por módulo. Kubernetes HPA
   responde automáticamente a picos de carga.}
  {Mayor complejidad operativa inicial: requiere herramientas de observabilidad
   distribuida (Jaeger, Prometheus), service mesh para comunicación segura entre
    servicios y equipos con experiencia en Kubernetes.}
    {Prueba de caos y carga: simular la caída del Motor de Alertas en producción
    controlada y verificar que los demás módulos mantienen operación. Criterio de
    aceptación: disponibilidad global $\geq$ 99,95\% y recuperación automática de pods
    en menos de 2 minutos.}

% ── ADR-002 ───────────────────────────────────────────────────────────────────
\adr{ADR-002 \textnormal{---} Edge Computing --- Gateways Locales en Subestaciones}
  {Aprobado}
  {Las subestaciones rurales de Colombia experimentan cortes de Internet frecuentes
   durante tormentas eléctricas, precisamente cuando más se necesita el monitoreo. La
   red eléctrica no puede esperar a que se restaure la conectividad para ejecutar
   acciones de protección.}
  {(1)~\textbf{Cloud-only sin edge}: toda la lógica en la nube; requiere conectividad
   constante --- inaceptable para zonas remotas.
   (2)~\textbf{Edge con lógica mínima (solo buffer)}: almacena pero no actúa; falla
   ante desconexión prolongada con evento crítico activo.
   (3)~\textbf{Edge con autonomía completa}: gateway evalúa reglas localmente, actúa
   y sincroniza después.}
  {Se instalan gateways ARM en cada subestación con InfluxDB local, motor de reglas
   CEP simplificado y capacidad de actuación local. Operan en modo autónomo durante
   desconexiones y sincronizan con la nube mediante MQTT QoS~2 al restaurar
   conectividad.}
  {Resiliencia ante cortes de Internet. Reducción de latencia para eventos críticos
   locales (sin round-trip a la nube). Protección física de la subestación incluso
   sin conectividad.}
  {Aumenta la superficie de ataque físico (dispositivos en campo sin vigilancia
   permanente). Requiere programa de actualización OTA robusto. El modelo de reglas
    en el edge debe mantenerse sincronizado con el modelo central.}
    {Ejecutar una prueba de desconexión regional de 24 horas. Criterio de aceptación:
    cero pérdida de eventos críticos, continuidad de alarmas locales y
    resincronización completa en menos de 15 minutos tras recuperar conectividad.}

% ── ADR-003 ───────────────────────────────────────────────────────────────────
\adr{ADR-003 \textnormal{---} Protocolo de Ingesta MQTT~v5 con Apache Kafka}
  {Aprobado}
  {La plataforma debe ingestar telemetría de 120.000 sensores simultáneos en redes
   celulares rurales de bajo ancho de banda (2G/3G en algunas zonas). La capa de
   ingesta es el punto de entrada de todo el sistema y no puede perder eventos
   críticos.}
  {(1)~\textbf{HTTP/REST}: alto overhead por cabeceras en cada mensaje; colapsa redes
   inestables con 120.000 dispositivos.
   (2)~\textbf{AMQP}: más robusto que HTTP pero mayor overhead que MQTT; pensado para
   mensajería empresarial, no IoT a escala.
   (3)~\textbf{MQTT~v5 + Kafka}: MQTT maneja el transporte ligero con QoS garantizado;
   Kafka proporciona el buffer de alta durabilidad y distribución a múltiples
   consumidores.}
  {MQTT~v5 como protocolo de transporte entre sensores/gateways y el broker. Apache
   Kafka como bus de eventos central con 24~particiones (1 por región de Colombia) y
   factor de replicación~3.}
  {Mínimo consumo de ancho de banda. QoS~2 garantiza entrega exactly-once para eventos
   críticos. Kafka permite múltiples consumidores independientes (Flink, ML, auditoría)
   del mismo stream sin acoplamiento.}
  {Requiere operar un clúster Kafka de alta disponibilidad (mínimo 3~brokers + KRaft).
    La combinación MQTT+Kafka agrega una capa de infraestructura adicional a monitorear
    y operar.}
    {Prueba de estrés de ingesta con 120.000 sensores simulados y picos del 60\%.
    Criterio de aceptación: throughput sostenido de al menos 50.000 eventos/s,
    escalamiento a 80.000 eventos/s en picos y lag de consumidores dentro del umbral
    operativo definido.}

% ── ADR-004 ───────────────────────────────────────────────────────────────────
\adr{ADR-004 \textnormal{---} Almacenamiento Multi-Tier (InfluxDB + Object Storage)}
  {Aprobado}
  {500~TB iniciales con crecimiento de 80~TB/año y obligación de retener datos 5~años.
   Los costos de mantener toda la telemetría en almacenamiento de alto rendimiento
   serían prohibitivos. Al mismo tiempo, el mantenimiento predictivo requiere acceso
   eficiente a años de históricos.}
  {(1)~\textbf{PostgreSQL + TimescaleDB}: buen soporte de series de tiempo pero no
   optimizado para el volumen de escritura requerido (50K eventos/s).
   (2)~\textbf{InfluxDB standalone}: óptimo para escritura pero costoso para retención
   a 5~años.
   (3)~\textbf{InfluxDB (hot) + Object Storage comprimido (cold)}: separa el acceso
   por temperatura de datos --- más eficiente en costo y rendimiento.}
  {Estrategia de tres niveles: InfluxDB Cluster para datos hot (90~días, acceso
   frecuente), Object Storage con compresión Parquet para datos warm (91~días --
   1~año) y cold (1--5~años). PostgreSQL para metadatos transaccionales y datos de
   usuarios.}
  {Reducción de costos de almacenamiento estimada en un 70\% respecto a mantener todo
   en InfluxDB. InfluxDB optimizado para escritura masiva de series de tiempo. Object
   Storage escala a petabytes sin rediseño arquitectónico.}
  {Mayor complejidad operativa: dos motores de bases de datos (InfluxDB + PostgreSQL)
    más gestión de Object Storage. Las consultas que cruzan datos hot y cold requieren
    federación mediante Apache Spark.}
    {Benchmark trimestral de costo y rendimiento por tier. Criterio de aceptación:
    reducción de costo total de almacenamiento frente a estrategia hot-only,
    cumplimiento de retención de 5 años y tiempos de consulta dentro de los SLA por
    tipo de dato (hot/warm/cold).}

% ── ADR-005 ───────────────────────────────────────────────────────────────────
\adr{ADR-005 \textnormal{---} Autenticación Mutua mTLS con Certificados X.509 por Dispositivo}
  {Aprobado}
  {La infraestructura eléctrica es objetivo de actores estatales y grupos criminales.
   Una autenticación débil puede permitir inyectar datos falsos que provoquen apagados
   innecesarios o que se ignoren anomalías reales.}
  {(1)~\textbf{Usuario/contraseña}: fácil de comprometer; escala muy mal para 120.000
   dispositivos.
   (2)~\textbf{API Keys compartidas}: si una clave se compromete, todos los dispositivos
   que la comparten quedan expuestos.
   (3)~\textbf{X.509 único por dispositivo con mTLS}: cada dispositivo tiene su
   certificado único; la revocación es quirúrgica; cumple NERC~CIP.}
  {Cada dispositivo físico (sensor, gateway) recibe un certificado X.509 único emitido
   por la PKI interna. La comunicación usa mTLS en todos los canales. La clave privada
   se almacena en el TPM~2.0 del hardware, no exportable por software.}
  {Revocación quirúrgica de dispositivos comprometidos sin afectar a los demás.
   Cumplimiento de NERC~CIP e ISO~27001 control A.10.1. La clave privada en TPM impide
   su extracción incluso con acceso físico al dispositivo.}
  {Requiere implementar y operar una PKI interna completa (CA raíz offline en HSM,
   CA intermedia online, OCSP responder). El aprovisionamiento inicial de 120.000
    certificados es operativamente complejo y debe automatizarse.}
    {Pruebas de seguridad continuas (red team + pruebas de intrusión) y escenarios de
    certificados revocados. Criterio de aceptación: 100\% de conexiones sin certificado
    válido bloqueadas, revocación efectiva en tiempo objetivo y trazabilidad completa
    en SIEM.}

% ── ADR-006 ───────────────────────────────────────────────────────────────────
\adr{ADR-006 \textnormal{---} Apache Flink como Motor de Stream Processing}
  {Aprobado}
  {El SLA de detección de anomalías en menos de 3~s exige un motor de procesamiento
   streaming de muy baja latencia capaz de manejar 50.000 eventos/s con semántica
   exactly-once y manejo de eventos fuera de orden por desincronización de reloj.}
  {(1)~\textbf{Apache Spark Streaming} (micro-batches): latencias de 100~ms -- 1~s por
   diseño; insuficiente para el SLA de 3~s.
   (2)~\textbf{Apache Storm}: baja latencia pero sin soporte nativo de estado y
   semántica exactly-once débil.
   (3)~\textbf{Apache Flink}: procesamiento de stream verdadero (event-driven),
   latencia sub-segundo, estado distribuido gestionado, soporte de event time con
   watermarks para corrección de eventos tardíos.}
  {Apache Flink como motor principal de stream processing. Se usa event time con
   watermarks de PTP para manejar desincronización de reloj. Las ventanas de tiempo
   son de 1~s, 5~s y 1~min para diferentes tipos de detección.}
  {Latencia de procesamiento sub-segundo. Manejo nativo de eventos tardíos mediante
   watermarks. Semántica exactly-once para garantizar que ninguna anomalía se detecte
   dos veces ni se pierda. Escalado horizontal automático con Kubernetes.}
  {Mayor curva de aprendizaje que Spark para el equipo. La gestión del estado
   distribuido requiere cuidado especial en actualizaciones de versión. El
    checkpointing frecuente genera overhead adicional de I/O.}
    {Pruebas de latencia y exactitud con eventos fuera de orden. Criterio de
    aceptación: latencia de detección p99,9 menor a 3 segundos, recuperación correcta
    desde checkpoints y semántica effectively-once en flujos críticos.}

% ── ADR-007 ───────────────────────────────────────────────────────────────────
\adr{ADR-007 \textnormal{---} Sincronización de Tiempo con PTP IEEE~1588}
  {Aprobado}
  {Los 120.000 sensores distribuidos en todo el país tienen relojes internos que se
   desincronizan con el tiempo, especialmente durante apagones y reinicios. La
   correlación correcta de eventos entre sensores de distintas regiones requiere
   timestamps precisos a nivel de microsegundos.}
  {(1)~\textbf{NTP}: precisión de milisegundos; insuficiente para correlacionar eventos
   eléctricos que ocurren en microsegundos.
   (2)~\textbf{GPS por sensor}: precisión sub-microsegundo pero costoso para 120.000
   dispositivos.
   (3)~\textbf{PTP IEEE~1588 con Grandmaster Clock por región}: precisión de
   microsegundos distribuida sin GPS en cada sensor.}
  {Despliegue de un PTP Grandmaster Clock por región sincronizado con GPS. Los
   gateways edge actúan como Boundary Clocks que distribuyen la señal de tiempo a los
   sensores locales. Apache Flink usa los timestamps PTP como event time con watermarks
   adaptativos.}
  {Correlación precisa de eventos de múltiples sensores con precisión de microsegundos.
   Eliminación de inconsistencias temporales en el análisis de apagones en cascada.
   Cumple el requisito de IEC~61850 para marcas de tiempo de protección.}
  {Requiere hardware de red compatible con PTP (switches con hardware timestamping) en
    las subestaciones. El coste inicial de los Grandmaster Clocks y su infraestructura
    de red asociada es elevado.}
    {Monitoreo continuo de offset y jitter por región. Criterio de aceptación:
    desincronización dentro del umbral de microsegundos definido para protección
    eléctrica y consistencia temporal en correlación de eventos multi-región.}

   % ══════════════════════════════════════════════════════════════════════════════
% SECCIÓN 9: ARQUITECTURA DE ALTO NIVEL (HLD)
% ══════════════════════════════════════════════════════════════════════════════
\section{Arquitectura de Alto Nivel (HLD)}

La arquitectura se organiza en seis zonas funcionales que reflejan el recorrido completo
de un evento desde el sensor físico hasta la decisión operativa del centro de control.
Las zonas se comunican de forma asíncrona a través de un bus de eventos central,
garantizando el desacoplamiento necesario para que cada componente escale de forma
independiente. Un plano transversal de seguridad vigila todas las zonas sin interrumpir
el flujo de datos.

\begin{center}
  \includegraphics[width=\textwidth]{HLD}
  \captionof{figure}{Arquitectura de Alto Nivel --- Plataforma IoT de Monitoreo Industrial para
  Infraestructura Energética}
  \label{fig:hld}
\end{center}

\subsection{Zonas de la Arquitectura}

\begin{longtable}{
  P{2.8cm}
  P{4.2cm}
  P{6.7cm}
}
\rowcolor{azulPUJ}
\multicolumn{1}{>{\raggedright\arraybackslash\color{white}\bfseries}p{2.8cm}}{Zona} &
\multicolumn{1}{>{\raggedright\arraybackslash\color{white}\bfseries}p{4.2cm}}{Componentes} &
\multicolumn{1}{>{\raggedright\arraybackslash\color{white}\bfseries}p{6.7cm}}{Responsabilidad} \\
\endfirsthead
\rowcolor{azulPUJ}
\multicolumn{1}{>{\raggedright\arraybackslash\color{white}\bfseries}p{2.8cm}}{Zona} &
\multicolumn{1}{>{\raggedright\arraybackslash\color{white}\bfseries}p{4.2cm}}{Componentes} &
\multicolumn{1}{>{\raggedright\arraybackslash\color{white}\bfseries}p{6.7cm}}{Responsabilidad} \\
\endhead

\rowcolor{azulClr}
\textbf{Zona de Campo} &
  Sensores IoT, Nodo Inteligente de Subestación, Centro de Control Heredado &
  Captura de telemetría física cada 5 segundos. El nodo inteligente procesa alarmas
  críticas localmente y almacena datos durante cortes de conectividad, garantizando
  autonomía operativa ante desconexiones prolongadas. \\

\rowcolor{azulClr}
\textbf{Transporte de Datos} &
  Receptor de Mensajes IoT, Bus de Eventos, Cola de Reintentos &
  Recepción masiva de eventos con entrega garantizada. El bus distribuye cada mensaje
  a todos los consumidores simultáneamente sin acoplamiento entre ellos. La cola de
  reintentos previene la pérdida de eventos críticos ante fallos transitorios. \\

\rowcolor{azulClr}
\textbf{Procesamiento Central} &
  Motor de Tiempo Real, Detector de Anomalías, Motor de Predicción IA,
  Sincronizador de Tiempo &
  Análisis de eventos en menos de 3 segundos. Detección de anomalías mediante reglas
  y modelos de inteligencia artificial. Anticipación de fallas futuras para mantenimiento
  predictivo. Corrección de desincronización de relojes entre regiones del país. \\

\rowcolor{azulClr}
\textbf{Memoria del Sistema} &
  Datos Recientes, Histórico Medio, Archivo Legal, Base de Datos Operacional &
  Almacenamiento escalonado por temperatura de datos: acceso inmediato para los
  últimos 90 días, consulta rápida hasta 1 año, y archivo regulatorio hasta 5 años.
  La base operacional gestiona usuarios, configuraciones y metadatos del sistema. \\

\rowcolor{azulClr}
\textbf{Servicios al Usuario} &
  Mapa Nacional, Centro de Notificaciones, Control Remoto Seguro,
  Exportación de Datos, Panel de Mantenimiento &
  Cinco puntos de salida adaptados a cada perfil de usuario: visualización geoespacial
  en tiempo real, alertas omnicanal, actuación remota con doble verificación de
  identidad, integración con sistemas externos y análisis predictivo para planificación
  de mantenimiento. \\

\rowcolor{azulClr}
\textbf{Seguridad Transversal} &
  Identidad Digital, Vigilancia Cibernética, Registro de Auditoría &
  Plano horizontal que vigila todas las zonas sin interrumpir el flujo de datos.
  Verifica la identidad de cada dispositivo y usuario, detecta comportamientos
  anómalos en tiempo real y mantiene un registro inmutable de todos los eventos
  del sistema para cumplimiento regulatorio. \\

\end{longtable}

\subsection{Flujo de una Alerta Crítica}

El siguiente diagrama describe el recorrido completo de un evento crítico desde
su origen físico hasta la actuación del operador, con los tiempos acumulados
en cada etapa para garantizar el SLA de 3~segundos.

\vspace{6pt}

\begin{longtable}{
  C{1.0cm}
  P{3.2cm}
  P{6.5cm}
  C{2.0cm}
}
\rowcolor{azulPUJ}
\multicolumn{1}{>{\centering\arraybackslash\color{white}\bfseries}p{1.0cm}}{Paso} &
\multicolumn{1}{>{\raggedright\arraybackslash\color{white}\bfseries}p{3.2cm}}{Componente} &
\multicolumn{1}{>{\raggedright\arraybackslash\color{white}\bfseries}p{6.5cm}}{Acción} &
\multicolumn{1}{>{\centering\arraybackslash\color{white}\bfseries}p{2.0cm}}{Latencia acum.} \\
\endfirsthead
\rowcolor{azulPUJ}
\multicolumn{1}{>{\centering\arraybackslash\color{white}\bfseries}p{1.0cm}}{Paso} &
\multicolumn{1}{>{\raggedright\arraybackslash\color{white}\bfseries}p{3.2cm}}{Componente} &
\multicolumn{1}{>{\raggedright\arraybackslash\color{white}\bfseries}p{6.5cm}}{Acción} &
\multicolumn{1}{>{\centering\arraybackslash\color{white}\bfseries}p{2.0cm}}{Latencia acum.} \\
\endhead

\cellcolor{azulClr}\textbf{1} & Sensor IoT &
  Captura lectura anómala y la publica con entrega garantizada y cifrado. &
  \cellcolor{verdeCl}0 ms \\
\cellcolor{azulClr}\textbf{2} & Nodo Inteligente &
  Recibe la señal, evalúa reglas críticas locales y reenvía hacia la nube. &
  \cellcolor{verdeCl}50--100 ms \\
\cellcolor{azulClr}\textbf{3} & Receptor + Bus de Eventos &
  Ingesta el mensaje y lo distribuye al procesamiento central. &
  \cellcolor{verdeCl}100--200 ms \\
\cellcolor{azulClr}\textbf{4} & Motor de Tiempo Real &
  Consume el evento, corrige el timestamp y evalúa ventanas temporales. &
  \cellcolor{verdeCl}300--800 ms \\
\cellcolor{azulClr}\textbf{5} & Detector de Anomalías &
  Clasifica el evento como Alerta Crítica Nivel~1 y lo publica. &
  \cellcolor{verdeCl}800 ms--1,5 s \\
\cellcolor{azulClr}\textbf{6} & Centro de Notificaciones &
  Envía SMS, correo y notificación push al operador responsable de la región. &
  \cellcolor{verdeCl}1,5--2,5 s \\
\cellcolor{azulClr}\textbf{7} & Mapa Nacional &
  Actualiza el marcador del sensor afectado en el panel geoespacial. &
  \cellcolor{naranCl}2,5--3,0 s \\

\end{longtable}

\subsection{Patrones Arquitectónicos Aplicados}

\begin{longtable}{
  P{3.5cm}
  P{4.5cm}
  P{5.7cm}
}
\rowcolor{azulPUJ}
\multicolumn{1}{>{\raggedright\arraybackslash\color{white}\bfseries}p{3.5cm}}{Patrón} &
\multicolumn{1}{>{\raggedright\arraybackslash\color{white}\bfseries}p{4.5cm}}{Aplicación en el sistema} &
\multicolumn{1}{>{\raggedright\arraybackslash\color{white}\bfseries}p{5.7cm}}{Beneficio principal} \\
\endfirsthead
\rowcolor{azulPUJ}
\multicolumn{1}{>{\raggedright\arraybackslash\color{white}\bfseries}p{3.5cm}}{Patrón} &
\multicolumn{1}{>{\raggedright\arraybackslash\color{white}\bfseries}p{4.5cm}}{Aplicación en el sistema} &
\multicolumn{1}{>{\raggedright\arraybackslash\color{white}\bfseries}p{5.7cm}}{Beneficio principal} \\
\endhead

\rowcolor{grisCl}
\textbf{Arquitectura Orientada a Eventos} &
  Bus central como único canal de comunicación entre zonas. &
  Desacoplamiento total; cada zona escala sin afectar a las demás. \\

\textbf{Separación Comando / Consulta} &
  Escrituras de telemetría separadas de las lecturas históricas. &
  Los análisis pesados no degradan el monitoreo en tiempo real. \\

\rowcolor{grisCl}
\textbf{Cortocircuito} &
  El nodo de subestación prioriza eventos críticos durante saturación. &
  Previene la propagación de fallos en cascada durante picos de carga. \\

\textbf{Estrangulador} &
  Integración gradual con centros de control heredados sin reemplazarlos. &
  Migración incremental sin riesgo de interrupción de sistemas existentes. \\

\rowcolor{grisCl}
\textbf{Mamparo} &
  Particionamiento del bus por región geográfica del país. &
  Un fallo en una región no afecta el monitoreo de las demás. \\

\textbf{Almacenar y Reenviar} &
  El nodo local guarda datos durante desconexiones y los envía al restaurarse. &
  Garantiza integridad de datos sin pérdida ante cortes de conectividad. \\

\end{longtable}

% ══════════════════════════════════════════════════════════════════════════════
% SECCIÓN 10: DIAGRAMA DE DESPLIEGUE
% ══════════════════════════════════════════════════════════════════════════════
\section{Diagrama de Despliegue}

La vista de despliegue describe dónde se ejecuta cada componente, cómo se conectan
las zonas OT/IT y qué mecanismos permiten continuidad operativa ante fallos
regionales.

\subsection{Vista de Despliegue Implementada}

\begin{center}
  \includegraphics[width=\textwidth]{despliegue.png}
  \captionof{figure}{Diagrama de despliegue de la plataforma IoT para monitoreo
  industrial de infraestructura energética.}
  \label{fig:despliegue}
\end{center}

\subsection{Lectura e Interpretación del Diagrama}

El diagrama \ref{fig:despliegue} representa una arquitectura distribuida en tres
planos: \textbf{campo (OT)}, \textbf{procesamiento en nube (IT)} y
\textbf{plataforma de datos históricos}. Esta separación permite aislar
responsabilidades operativas y escalar cada plano de forma independiente según su
patrón de carga.

En el plano OT, los gateways edge reciben telemetría de sensores y ejecutan lógica
crítica local para mantener operación segura cuando la conectividad se degrada. En
el plano IT, los clústeres regionales activos-activos absorben la ingesta masiva,
procesan eventos de streaming y publican alertas en tiempo real. Finalmente, el
plano de datos implementa retención por niveles para equilibrar rendimiento de
consulta y costo total de almacenamiento.

La integración entre regiones no solo responde a disponibilidad, sino también a
objetivos de continuidad: ante caída de una región, la segunda mantiene operación
sin pérdida de funciones críticas. Este comportamiento debe leerse junto con los
SLA de RTO/RPO y con los ADR de tolerancia a fallos y arquitectura edge.

\subsection{Decisiones de Despliegue y Justificación Técnica}

\begin{longtable}{P{3.7cm} P{4.0cm} P{5.9cm}}
\rowcolor{azulPUJ}
\multicolumn{1}{>{\raggedright\arraybackslash\color{white}\bfseries}p{3.7cm}}{Decisión} &
\multicolumn{1}{>{\raggedright\arraybackslash\color{white}\bfseries}p{4.0cm}}{Necesidad del caso} &
\multicolumn{1}{>{\raggedright\arraybackslash\color{white}\bfseries}p{5.9cm}}{Efecto esperado} \\
\endfirsthead
\rowcolor{azulPUJ}
\multicolumn{1}{>{\raggedright\arraybackslash\color{white}\bfseries}p{3.7cm}}{Decisión} &
\multicolumn{1}{>{\raggedright\arraybackslash\color{white}\bfseries}p{4.0cm}}{Necesidad del caso} &
\multicolumn{1}{>{\raggedright\arraybackslash\color{white}\bfseries}p{5.9cm}}{Efecto esperado} \\
\endhead

\rowcolor{grisCl}
Edge en subestaciones &
  Pérdida de conectividad en nodos remotos y eventos climáticos extremos &
  Continuidad local, protección operativa y envío diferido sin pérdida de eventos críticos. \\

Activa-activa multi-región &
  Disponibilidad anual de 99,95\% y tolerancia a fallos regionales &
  Alta disponibilidad, failover automático y reducción de interrupciones de servicio. \\

\rowcolor{grisCl}
Bus de eventos desacoplado &
  Throughput masivo y consumidores heterogéneos (alertas, analítica, auditoría) &
  Escalado horizontal de ingesta y procesamiento sin acoplamiento rígido entre módulos. \\

Almacenamiento por niveles (hot/warm/cold) &
  Retención regulatoria de 5 años con crecimiento anual de volumen &
  Balance entre costo y rendimiento, garantizando consulta rápida reciente y archivo legal. \\

\rowcolor{grisCl}
Seguridad transversal (PKI + SIEM) &
  Riesgo de ciberataques a infraestructura crítica &
  Autenticación fuerte, detección temprana de incidentes y trazabilidad auditada de extremo a extremo. \\

\end{longtable}

\subsection{Nodos de Despliegue y Responsabilidades}

\begin{longtable}{P{3.8cm} P{3.8cm} P{2.6cm} P{3.5cm}}
\rowcolor{azulPUJ}
\multicolumn{1}{>{\raggedright\arraybackslash\color{white}\bfseries}p{3.8cm}}{Nodo} &
\multicolumn{1}{>{\raggedright\arraybackslash\color{white}\bfseries}p{3.8cm}}{Componentes} &
\multicolumn{1}{>{\centering\arraybackslash\color{white}\bfseries}p{2.6cm}}{Zona} &
\multicolumn{1}{>{\raggedright\arraybackslash\color{white}\bfseries}p{3.5cm}}{Objetivo de diseño} \\
\endfirsthead
\rowcolor{azulPUJ}
\multicolumn{1}{>{\raggedright\arraybackslash\color{white}\bfseries}p{3.8cm}}{Nodo} &
\multicolumn{1}{>{\raggedright\arraybackslash\color{white}\bfseries}p{3.8cm}}{Componentes} &
\multicolumn{1}{>{\centering\arraybackslash\color{white}\bfseries}p{2.6cm}}{Zona} &
\multicolumn{1}{>{\raggedright\arraybackslash\color{white}\bfseries}p{3.5cm}}{Objetivo de diseño} \\
\endhead

\rowcolor{grisCl}
Gateway de Subestacion &
  Buffer local, motor de reglas criticas, cliente MQTT con QoS~2 &
  OT/Edge &
  Operacion autonoma ante desconexion y cero perdida de eventos criticos. \\

Cluster Kubernetes Region A &
  API Gateway, Broker MQTT, Kafka, Flink, servicios de negocio &
  IT/Cloud &
  Baja latencia y escalado horizontal para alta carga de eventos. \\

\rowcolor{grisCl}
Cluster Kubernetes Region B &
  Misma capacidad de Region A en modo activo-activo &
  IT/Cloud &
  Tolerancia a fallos regionales con continuidad de servicio. \\

Plataforma de Datos Historicos &
  InfluxDB hot, Data Lake warm/cold, almacenamiento WORM &
  IT/Data &
  Retencion de 5 anos optimizando costo y rendimiento de consulta. \\

\rowcolor{grisCl}
Plataforma de Seguridad y Observabilidad &
  PKI X.509, SIEM, WAF/IDS, trazas y metricas &
  Transversal &
  Visibilidad operativa, deteccion temprana de incidentes y cumplimiento. \\

\end{longtable}

% ══════════════════════════════════════════════════════════════════════════════
% SECCIÓN 11: ENTREGABLES
% ══════════════════════════════════════════════════════════════════════════════
\section{Entregables}

\subsection{Documento}

El documento final cumple los elementos solicitados para el caso de estudio:

\begin{longtable}{P{9.5cm} C{4.2cm}}
\rowcolor{azulPUJ}
\multicolumn{1}{>{\raggedright\arraybackslash\color{white}\bfseries}p{9.5cm}}{Elemento requerido} &
\multicolumn{1}{>{\centering\arraybackslash\color{white}\bfseries}p{4.2cm}}{Estado} \\
\endfirsthead
\rowcolor{azulPUJ}
\multicolumn{1}{>{\raggedright\arraybackslash\color{white}\bfseries}p{9.5cm}}{Elemento requerido} &
\multicolumn{1}{>{\centering\arraybackslash\color{white}\bfseries}p{4.2cm}}{Estado} \\
\endhead

\rowcolor{grisCl}
Requerimientos (proyecto, dominio y producto) & Completado \\
Usuarios & Completado \\
\rowcolor{grisCl}
Matriz Modulos Funcionales vs Usuarios & Completado \\
Restricciones y riesgos & Completado \\
\rowcolor{grisCl}
Atributos de calidad ISO~25010 (descripcion, prioridad y clasificacion) & Completado \\
Escenarios de calidad (minimo 1 por subcategoria priorizada) & Completado \\
\rowcolor{grisCl}
Acuerdos de atributos de calidad (SLA) & Completado \\
ADRs de la arquitectura propuesta & Completado \\
\rowcolor{grisCl}
Arquitectura de Alto Nivel (HLD) & Completado \\
Diagrama de despliegue & Completado \\

\end{longtable}


\subsection{Lecciones Aprendidas y Conclusiones}

\textbf{Lecciones aprendidas.}
La estrategia orientada a eventos y despliegue multi-region permite responder al
throughput exigido sin sacrificar resiliencia. La incorporacion de capacidades
edge no es opcional en entornos con conectividad intermitente. La seguridad debe
disenarse de forma transversal desde el inicio para evitar retrabajos costosos.

\textbf{Conclusiones.}
La arquitectura propuesta satisface los requisitos criticos de latencia,
disponibilidad y retencion regulatoria, manteniendo control de costos mediante
almacenamiento por niveles. Los ADRs y SLAs definidos ofrecen trazabilidad tecnica
y criterios verificables para evolucionar la plataforma en produccion.

\end{document}